% GENERATED FILE. Edit canonical lessons, then run npm run generate:books.
% canonical-source-hash: fnv1a64:0b50ea898de91fa9
% canonical-lessons: GE-C507-halbzeit, GE-C507-elfmeter, GE-C507-verlaengerung, GE-C507-unentschieden, GE-C507-abseits

\chapter{Halbzeit und Elfmeter}
\label{ch:ge-halbzeit-und-elfmeter}
\hlchaptermodality{507}

\begin{chapteropening}
\textbf{By the end of this chapter:} \emph{I can talk about the stages of a match.}

Complete the last lesson of chapter 507: I can talk about the stages of a match.
\end{chapteropening}

\section[die Halbzeit]{die Halbzeit — half-time; half}

\label{lesson:GE-C507-halbzeit}



\begin{quote}
Before the new one: say the German for referee, then the German for to whistle.
\end{quote}

\subsection*{You'll want to know: die Halbzeit}
\textbf{die Halbzeit} — \textquotedblleft{}half-time; half\textquotedblright{}.

\begin{quote}
In der Halbzeit trinken die Spieler Wasser.
\end{quote}

\begin{quote}
\emph{At half-time the players drink water.}
\end{quote}

Plural: \emph{die Halbzeiten}. \emph{Die erste Halbzeit} — the first half.

\begin{cousinweb}
\emph{Halbzeit} is built on \emph{halb} \textquotedblleft{}half\textquotedblright{}, from Proto-Germanic \emph{halbaz}, and \emph{Zeit} \textquotedblleft{}time\textquotedblright{}, inherited from Old High German \emph{zīt}, from Proto-Germanic \emph{tīdiz} (EtymDB).
\end{cousinweb}

\begin{grammarlens}[title={What you've built}]
One more word for telling, explaining and describing at B1.
\end{grammarlens}

\subsection*{Your turn}
\begin{itemize}
\raggedright
  \item \emph{Say it:} \emph{die Halbzeit}
  \item \emph{Say it:} \emph{die Halbzeit}, once more
  \item \emph{Say it:} a sentence of your own with \emph{die Halbzeit}
  \item \emph{Recall:} say the German for referee, then the German for to whistle, then say \emph{die Halbzeit} again
\end{itemize}

\begin{tcolorbox}[breakable,colback=teal!4,colframe=teal!35!black,title={Before you move on}]
What does die Halbzeit mean? (\textquotedblleft{}Half-time; half\textquotedblright{}.) Say it once more.
\end{tcolorbox}
\section[der Elfmeter]{der Elfmeter — penalty (kick)}

\label{lesson:GE-C507-elfmeter}



\begin{quote}
Before the new one: say the German for to whistle, then the German for half-time.
\end{quote}

\subsection*{You'll want to know: der Elfmeter}
\textbf{der Elfmeter} — \textquotedblleft{}penalty (kick)\textquotedblright{}.

\begin{quote}
In der letzten Minute gab es einen Elfmeter.
\end{quote}

\begin{quote}
\emph{In the last minute there was a penalty.}
\end{quote}

Plural: \emph{die Elfmeter}. Named after the eleven metres to the goal.

\begin{cousinweb}
\emph{Elfmeter} is built on \emph{elf} \textquotedblleft{}eleven\textquotedblright{} and \emph{Meter}, which was borrowed from French \emph{mètre}, from Greek \emph{métron} \textquotedblleft{}measure\textquotedblright{} (EtymDB).
\end{cousinweb}

\begin{grammarlens}[title={What you've built}]
One more word for telling, explaining and describing at B1.
\end{grammarlens}

\subsection*{Your turn}
\begin{itemize}
\raggedright
  \item \emph{Say it:} \emph{der Elfmeter}
  \item \emph{Say it:} \emph{der Elfmeter}, once more
  \item \emph{Say it:} a sentence of your own with \emph{der Elfmeter}
  \item \emph{Recall:} say the German for to whistle, then the German for half-time, then say \emph{der Elfmeter} again
\end{itemize}

\begin{tcolorbox}[breakable,colback=teal!4,colframe=teal!35!black,title={Before you move on}]
What does der Elfmeter mean? (\textquotedblleft{}Penalty (kick)\textquotedblright{}.) Say it once more.
\end{tcolorbox}
\section[die Verlängerung]{die Verlängerung — extra time; extension}

\label{lesson:GE-C507-verlaengerung}



\begin{quote}
Before the new one: say the German for half-time, then the German for penalty.
\end{quote}

\subsection*{You'll want to know: die Verlängerung}
\textbf{die Verlängerung} — \textquotedblleft{}extra time; extension\textquotedblright{}.

\begin{quote}
Nach der Verlängerung stand es immer noch eins zu eins.
\end{quote}

\begin{quote}
\emph{After extra time it was still one all.}
\end{quote}

Plural: \emph{die Verlängerungen}. Also the extension of a contract or a cable.

\begin{cousinweb}
\emph{Verlängerung} is built on \emph{verlängern} \textquotedblleft{}to lengthen\textquotedblright{}, from \emph{lang} \textquotedblleft{}long\textquotedblright{}, inherited from Old High German \emph{lang}, from Proto-Germanic \emph{langaz} (EtymDB).
\end{cousinweb}

\begin{grammarlens}[title={What you've built}]
One more word for telling, explaining and describing at B1.
\end{grammarlens}

\subsection*{Your turn}
\begin{itemize}
\raggedright
  \item \emph{Say it:} \emph{die Verlängerung}
  \item \emph{Say it:} \emph{die Verlängerung}, once more
  \item \emph{Say it:} a sentence of your own with \emph{die Verlängerung}
  \item \emph{Recall:} say the German for half-time, then the German for penalty, then say \emph{die Verlängerung} again
\end{itemize}

\begin{tcolorbox}[breakable,colback=teal!4,colframe=teal!35!black,title={Before you move on}]
What does die Verlängerung mean? (\textquotedblleft{}Extra time; extension\textquotedblright{}.) Say it once more.
\end{tcolorbox}
\section[das Unentschieden]{das Unentschieden — tied game, level score}

\label{lesson:GE-C507-unentschieden}



\begin{quote}
Before the new one: say the German for penalty, then the German for extra time.
\end{quote}

\subsection*{You'll want to know: das Unentschieden}
\textbf{das Unentschieden} — \textquotedblleft{}tied game, level score\textquotedblright{}.

\begin{quote}
Das Spiel endete mit einem Unentschieden.
\end{quote}

\begin{quote}
\emph{The game ended level.}
\end{quote}

Plural: \emph{die Unentschieden}. As an adjective: \emph{Das Spiel ging unentschieden aus.}

\begin{cousinweb}
\emph{Unentschieden} is built on \emph{un-} \textquotedblleft{}not\textquotedblright{} and \emph{entschieden} \textquotedblleft{}decided\textquotedblright{}, the participle of \emph{entscheiden} \textquotedblleft{}to decide\textquotedblright{} (EtymDB lists the verb): a game that nobody decided.
\end{cousinweb}

\begin{grammarlens}[title={What you've built}]
One more word for telling, explaining and describing at B1.
\end{grammarlens}

\subsection*{Your turn}
\begin{itemize}
\raggedright
  \item \emph{Say it:} \emph{das Unentschieden}
  \item \emph{Say it:} \emph{das Unentschieden}, once more
  \item \emph{Say it:} a sentence of your own with \emph{das Unentschieden}
  \item \emph{Recall:} say the German for penalty, then the German for extra time, then say \emph{das Unentschieden} again
\end{itemize}

\begin{tcolorbox}[breakable,colback=teal!4,colframe=teal!35!black,title={Before you move on}]
What does das Unentschieden mean? (\textquotedblleft{}Tied game, level score\textquotedblright{}.) Say it once more.
\end{tcolorbox}
\section[das Abseits]{das Abseits — offside}

\label{lesson:GE-C507-abseits}



\begin{quote}
Before the new one: say the German for extra time, then the German for tied game.
\end{quote}

\subsection*{You'll want to know: das Abseits}
\textbf{das Abseits} — \textquotedblleft{}offside\textquotedblright{}.

\begin{quote}
Das Tor zählte nicht, weil der Spieler im Abseits stand.
\end{quote}

\begin{quote}
\emph{The goal didn't count because the player was offside.}
\end{quote}

\emph{Im Abseits stehen} — to be offside; also to be left out.

\begin{cousinweb}
\emph{Abseits} is built on \emph{ab} \textquotedblleft{}off\textquotedblright{} and \emph{seits} \textquotedblleft{}on the side\textquotedblright{}; EtymDB links the football sense to English \emph{offside}, which it renders.
\end{cousinweb}

\begin{grammarlens}[title={What you've built}]
One more word for telling, explaining and describing at B1.
\end{grammarlens}

\subsection*{Your turn}
\begin{itemize}
\raggedright
  \item \emph{Say it:} \emph{das Abseits}
  \item \emph{Say it:} \emph{das Abseits}, once more
  \item \emph{Say it:} a sentence of your own with \emph{das Abseits}
  \item \emph{Recall:} say the German for extra time, then the German for tied game, then say \emph{das Abseits} again
\end{itemize}

\begin{tcolorbox}[breakable,colback=teal!4,colframe=teal!35!black,title={Before you move on}]
What does das Abseits mean? (\textquotedblleft{}Offside\textquotedblright{}.) Say it once more.
\end{tcolorbox}
